% !PRINTED 109
\weberpage{109}

Denn setzen wir
\begin{equation}
 f(x)=x^n-a,
\tag{1}
\end{equation}
so ist $f(x)$ nach §. 31, VII. für ein hinlänglich grosses positives $x$ positiv, und für $x=0$ negativ; also giebt es einen positiven Werth von $x$, den man mit $\sqrt[n]{a}$ bezeichnet, für den $f(x)$ verschwindet. Dass es aber nur einen solchen Werth geben kann, folgt daraus, dass $x^n$, so lange $x$ positiv bleibt, mit $x$ fortwährend wächst.

Ist $n$ ungerade, so existirt auch für negative $a$ eine und nur eine negative Wurzel von (1), was ebenso bewiesen wird.

\medskip
\noindent X. \gesperrt{Eine Gleichung ungeraden Grades mit reellen Coëfficienten hat immer wenigstens eine reelle Wurzel.}
\medskip

Denn nach Satz VII. kann $f(x)$, wenn der Grad ungerade ist, sowohl positiv als negativ werden und $f(x)=0$ hat also nach VIII. eine Wurzel.

Wir beweisen endlich noch

\medskip
\noindent XI. \gesperrt{Eine reine Gleichung hat immer eine Wurzel.}
\medskip

Unter einer reinen Gleichung verstehen wir eine Gleichung von der Form
\begin{equation}
 x^n-a=0,
\tag{2}
\end{equation}
wo $a$ eine beliebige complexe Grösse ist. Dass diese Gleichung für ein reelles positives $a$ immer eine Wurzel hat, ist oben schon bewiesen, und für $a=0$ wird sie offenbar durch $x=0$ befriedigt.

Der allgemeine Beweis kann in zwei Theile zerlegt werden; es genügt nämlich, wenn wir zunächst $n=2$, dann $n$ ungerade voraussetzen. Denn existirt die Grösse $\sqrt[n]{a}$ und $\sqrt[2]{a}$ oder $\sqrt{a}$ für jedes $a$, so ist auch die Existenz von
\[
 \sqrt[2n]{a}=\sqrt[2]{\sqrt[n]{a}},\qquad \sqrt[4n]{a}=\sqrt[2]{\sqrt[2n]{a}}\ldots,
\]
also von $\sqrt[n]{a}$ für jedes beliebige $n$ sichergestellt. Es würde sogar genügen, die Existenz von $\sqrt[n]{a}$ für den Fall zu beweisen, dass $n$ eine Primzahl ist. Daraus ist aber zunächst kein besonderer Vortheil zu ziehen.

Wir bezeichnen die complexe Grösse $a$ mit $b+ci$ und setzen, da $x$ im Allgemeinen auch complex sein wird,
\[
 x=y+iz.
\]
